% Chapter 23: Conversation and Everyday English: Talking with Confidence

\chapter{Conversation and Everyday English: Talking with Confidence}

\section{Pedagogical Purpose}
Learners have now built a strong grammatical foundation and can construct accurate sentences of many kinds. This chapter turns that knowledge outward, into real spoken communication — the everyday language of greeting people, making requests, asking questions, responding politely, and asking for clarification when something is unclear. It is less about a single grammar rule and more about using the grammar already learned appropriately and politely.

\begin{learningobjectives}
The learner can:
\begin{itemize}
    \item use appropriate greetings for different times of day and situations.
    \item make polite requests using suitable phrases.
    \item ask questions to get information or clarification.
    \item respond appropriately to greetings, requests, and questions.
    \item take part in a short role play using natural, polite classroom English.
\end{itemize}
\end{learningobjectives}

\section{Prerequisite Check}
\begin{quickcheck}
Which of these is a polite request, and which is a command?

\begin{enumerate}
    \item Pass the salt.
    \item Could you please pass the salt?
\end{enumerate}
\end{quickcheck}

\begin{rememberbox}
You already know how to form questions and commands (Chapter 2) and how to report what someone said (Chapter 22). Now you will use these skills to communicate politely and clearly in everyday situations.
\end{rememberbox}

\section{Chapter Opener}
Rohan is practising for a class role play, but Ms. Sharma notices something.

\textbf{Rohan:} \textit{(to an imaginary shopkeeper)} "Give me a pencil!"

\textbf{Ms. Sharma:} "Rohan, that would work at home with your little brother, perhaps! But if you walked into a shop and said that, how do you think the shopkeeper might feel?"

\textbf{Rohan:} \textit{(laughing)} "Oh... maybe a bit surprised! I should ask more politely."

\textbf{Maya:} "You could say, 'Could I have a pencil, please?'"

\textbf{Ms. Sharma:} "Exactly, Maya! The words we choose change how our message sounds, even if the basic idea is the same. Let's explore how to speak politely and clearly in different everyday situations."

\section{Language Experience}
\textbf{Two Ways of Asking}

\begin{tcolorbox}[colback=gray!5, colframe=gray!50, sharp corners, boxsep=0pt, left=5pt, right=5pt, top=5pt, bottom=5pt, breakable]
\textbf{Version A:} "Give me that book."
\textbf{Version B:} "Could I please borrow that book?"
\end{tcolorbox}

\textbf{A Short Conversation}

\begin{tcolorbox}[colback=gray!5, colframe=gray!50, sharp corners, boxsep=0pt, left=5pt, right=5pt, top=5pt, bottom=5pt, breakable]
\textbf{Maya:} Good morning! How are you today?
\textbf{Rohan:} I'm well, thank you. And you?
\textbf{Maya:} I'm fine, thanks. Sorry, could you repeat that? I didn't quite hear you.
\textbf{Rohan:} Of course --- I said, "How are you?"
\end{tcolorbox}

\section{Look and Notice}
\begin{thinkbox}
\begin{enumerate}
    \item Which version of the request (A or B) sounds more polite? Why?
    \item In the conversation, how does Maya greet Rohan?
    \item What phrase does Maya use when she doesn't understand something?
    \item How does Rohan respond when Maya asks him to repeat himself?
\end{enumerate}
\end{thinkbox}

\section{Discover / Explicit Teaching}
Everyday spoken English uses certain useful phrases for common purposes:

\begin{table}[h!]
\centering
\begin{tabular}{|l|l|}
\hline
\textbf{Purpose} & \textbf{Useful Phrases} \\
\hline
Greeting & "Good morning/afternoon/evening," "Hello," "Hi, how are you?" \\
Making a polite request & "Could you please...", "May I...", "Would you mind..." \\
Asking a question & "What...", "Where...", "Why...", "Can you tell me...?" \\
Responding politely & "Yes, of course," "I'm sorry, but...", "Thank you," "You're welcome." \\
Asking for clarification & "Sorry, could you repeat that?", "What do you mean?", "Could you explain that again?" \\
\hline
\end{tabular}
\end{table}

\section{Grammar Word}
\begin{grammarword}{REGISTER}
\textbf{Register} means choosing words and a tone that suit the situation and the person you are speaking to — for example, speaking more formally to a stranger or a teacher, and more casually to a close friend.

\textit{Examples:}
\begin{itemize}
    \item To a teacher: "Excuse me, may I ask a question?"
    \item To a close friend: "Hey, quick question!"
\end{itemize}

\textit{Non-example:} Using very casual language with a shopkeeper you have never met ("Oi, give me that!") does not match the situation and can sound rude, even if it is not meant that way.

\textit{Quick Check:} Which sounds more appropriate when speaking to your school principal: "What's up?" or "Good morning, may I speak with you?"
\end{grammarword}

\section{Core Explanation}

\subsection*{Greeting appropriately}
``Good morning, Ms. Sharma.'' (formal, to a teacher) vs. ``Hi!'' (casual, to a friend). \textit{Check:} Does my greeting match the time of day and the person I am speaking to?

\subsection*{Making polite requests}
``Could you please pass the ruler?'' vs. ``Give me the ruler.'' (blunt, less polite). \textit{Check:} Have I used a polite phrase like "could," "please," or "may I," instead of only a command?

\subsection*{Asking for clarification}
``Sorry, could you say that again, please?'' vs. Staying silent and guessing what was said can lead to misunderstandings. \textit{Check:} If I did not understand something, have I asked politely instead of guessing?

\subsection*{Responding appropriately}
``Thank you for asking! I'm doing well.'' (response to a greeting) vs. Ignoring a greeting or question is impolite, even if unintentional. \textit{Check:} Have I responded to what the other person actually said or asked?

\section{Worked Example}
\begin{examplebox}
\textbf{Situation:} You want to ask your teacher if you can leave the classroom to get some water.

\textit{Step 1:} Think about who you are speaking to — a teacher, so a polite, slightly formal register is appropriate.
\textit{Step 2:} Choose a polite request phrase — "May I..." or "Could I..." both work well for asking permission.
\textit{Step 3:} State the request clearly and add "please" for extra politeness.

\textbf{Answer:} \textbf{Excuse me, may I please go and get some water?}
\end{examplebox}

\section{Guided Practice}
\begin{activitybox}{Activity A: Identification}
Say whether each phrase is a greeting, a request, or a way of asking for clarification.
\begin{enumerate}
    \item "Good morning!"
    \item "Could you please pass the book?"
    \item "Sorry, could you repeat that?"
\end{enumerate}
\end{activitybox}

\begin{activitybox}{Activity B: Completion}
Make each blunt sentence more polite.
\begin{enumerate}
    \item "Give me a pencil." $\rightarrow$ \_\_\_
    \item "Tell me the time." $\rightarrow$ \_\_\_
\end{enumerate}
\end{activitybox}

\section{Independent Practice}
\begin{activitybox}{Independent Practice}
\begin{enumerate}
    \item Write a short greeting you would use to meet your school principal for the first time.
    \item Write a polite request asking a classmate to lend you their eraser.
    \item Write a sentence you could use if you did not hear what someone said clearly.
    \item \textit{Reasoning:} Why might the same request sound polite to one person and rude to another, depending on how it is said?
\end{enumerate}
\end{activitybox}

\section{Common Confusion}
\begin{warningbox}
Learners sometimes think that being polite means using very long or complicated sentences, when in fact even a short phrase can be polite if it includes the right words ("please," "could," "may I").

\textit{How to check:} Ask, "Have I included at least one polite word or phrase, even if my sentence is short?"
\end{warningbox}

\section{Writing Connection}
Write a short dialogue (6-8 lines) between two characters in a shop, a school office, or a library. Include at least one greeting, one polite request, and one polite response.

\section{Summary}
\begin{summarybox}
\begin{itemize}
    \item Choose greetings that match the time of day and the situation.
    \item Use polite phrases like "could you," "may I," and "please" for requests.
    \item Ask questions clearly, and ask for clarification politely if you do not understand.
    \item Respond appropriately to greetings, requests, and questions.
    \item \textbf{Register} means adjusting your language to suit the person and situation.
\end{itemize}
\end{summarybox}

\section{Key Terms}
\begin{table}[h!]
\centering
\begin{tabular}{|l|l|}
\hline
\textbf{Term} & \textbf{Simple Meaning} \\
\hline
Register & Choosing words and tone to suit the situation \\
Greeting & A polite way of meeting someone \\
Polite request & Asking for something in a courteous way \\
Clarification & Asking someone to explain or repeat something unclear \\
Role play & Acting out a conversation to practise real-life language \\
\hline
\end{tabular}
\end{table}

\begin{selfcheck}
I can:
\begin{itemize}
    \item use appropriate greetings for different situations.
    \item make a polite request.
    \item ask a question to get information or clarification.
    \item respond appropriately in a conversation.
    \item take part in a short role play using natural, polite English.
\end{itemize}
\end{selfcheck}